\section{总结与延伸}

本文梳理了使用 OpenAI Codex 进行数据分析与报告交付的完整工作流，核心要点如下：

\begin{enumerate}
    \item \textbf{问题先行}：先确定一个具体、可验证的分析问题，再动手写代码
    \item \textbf{规范先行}：通过 \texttt{AGENTS.md} 定义项目规范，让 Codex 在一致的约束下工作
    \item \textbf{数据先行}：先盘点数据（inventory），再清洗合并，``先 Profile 再 Merge''
    \item \textbf{隔离探索}：用 Worktree 将不同方向的 EDA 物理隔离，保持 diff 可审查
    \item \textbf{基线优先}：从可解释模型开始，弱结果本身也是诊断信号
    \item \textbf{交付导向}：最终输出的不是 notebook，而是受众可直接使用的 artifact
    \item \textbf{技能复用}：将稳定的工作流封装为 Skill，提升长期效率
\end{enumerate}

\subsection{拓展阅读}

\begin{itemize}
    \item OpenAI Codex 官方文档：\url{https://developers.openai.com/codex/}
    \item Codex Skills 文档：\url{https://developers.openai.com/codex/skills/}
    \item 《R for Data Science》数据分析循环框架：\url{https://r4ds.had.co.nz/}
    \item Codex Worktree 使用指南：\url{https://developers.openai.com/codex/worktrees/}
\end{itemize}

%% --- Fin del contenido --- %%

\end{document}
